\section{Заключение}

В работе были изучены актуальные методы верификации, исследованы их преимущества и недостатки. На примере кэша второго уровня графического процессора общшегно назначения были исследованы такие характеристики как качество верификации, скорость нахождения ошибок и полнота. На основе полученных результатов были сформированы практические рекомендации о применении соответствующих методов. 

Особое внимание было уделено формальной верификации, как относительно новому и и неосвещенному методу. Были изучены применимость этого метода по отношению к размеру проверяемой системы.

В будущем развитие данной работы может представлять ...